\chapter{Минимальные данные производного отображения момента}

После закрытия обычного и скрученного спектральных следов остаётся
производная симплектическая ветвь. Её наиболее экономный кандидат ---
предпроективная алгебра удвоенного трёхвершинного колчана. Проверяется не
само известное тождество отображения момента, а более сильное утверждение:
следует ли требуемая ориентация из уже принятой вещественной геометрии KO6
без нового выбора поляризации, калибровочной группы или относительного
следа.

\section{Удвоенный колчан и точное локальное соотношение}

Пусть исходный остов имеет вершины \(0,G,2\) и стрелки
\begin{equation}
 0\xrightarrow{x}G\xrightarrow{y}2.
\end{equation}
В удвоенном колчане присутствуют формальные обратные стрелки \(x^*,y^*\).
Предпроективные соотношения по вершинам имеют вид
\begin{align}
 r_0&=-x^*x,\\
 r_G&=xx^*-y^*y,\\
 r_2&=yy^*.
 \label{eq:v5-derived-preprojective-relations}
\end{align}
При представлении
\begin{equation}
 \pi(x)=X,\qquad \pi(x^*)=X^\dagger,\qquad
 \pi(y)=Y,\qquad \pi(y^*)=Y^\dagger
\end{equation}
средняя компонента точно равна
\begin{equation}
 \boxed{\pi(r_G)=XX^\dagger-Y^\dagger Y.}
 \label{eq:v5-derived-exact-middle-relation}
\end{equation}
Это совпадает с блоком коммутатора Ходжа
\begin{equation}
 [d,d^\dagger]\big|_{V_G}=XX^\dagger-Y^\dagger Y,\qquad
 d=\begin{pmatrix}0&0&0\\X&0&0\\0&Y&0\end{pmatrix}.
\end{equation}
Следовательно, предпроективный язык без подбора коэффициентов воспроизводит
правильное алгебраическое выражение. Это содержательный проход, но пока не
доказательство его родительского происхождения.

\section{Две различные независимости от ориентации}

Абстрактная предпроективная алгебра не зависит от первоначального выбора
направлений рёбер. Однако соответствующий изоморфизм использует знак.
Если первое ребро обращается и новая исходная стрелка обозначается
\(b=x^*\), то для сохранения соотношения требуется
\begin{equation}
 b^*=-x.
 \label{eq:v5-derived-formal-orientation-map}
\end{equation}
Действительно, вклад средней вершины становится
\begin{equation}
 -b^*b=-(-x)x^*=xx^*.
\end{equation}
Но в положительном гильбертовом представлении обратная стрелка должна быть
сопряжённой:
\begin{equation}
 \pi(b^*)=\pi(b)^\dagger=X,
\end{equation}
тогда как формула \eqref{eq:v5-derived-formal-orientation-map} требует
\(\pi(b^*)=-X\). Поэтому алгебраический изоморфизм не сохраняет физическую
инволюцию. При сохранении инволюции вклад меняет знак:
\begin{equation}
 -b^*b=-XX^\dagger.
\end{equation}
Препятствие не устраняется фазовым переопределением. В общем случае
\(b=\alpha x^*\), \(b^*=\beta x\). Сохранение исходного вклада требует
\(\alpha\beta=-1\), тогда как положительная инволюция даёт
\(\beta=\overline{\alpha}\) и, следовательно,
\(\alpha\beta=|\alpha|^2\geq0\). Эти условия несовместимы.
Та же картина возникает на втором ребре. Одновременное обращение всех
стрелок меняет \(r_G\mapsto-r_G\) и сохраняет \(\Tr r_G^2\), но независимое
обращение одного ребра превращает цепь в сток или источник и меняет
функционал. Это в точности неединственность, ранее найденная для оператора
высоты.

\begin{theorem}[Несовместимость двух требований каноничности]
Независимость предпроективного соотношения от обращения отдельного ребра и
условие «формальная обратная стрелка представляется эрмитово сопряжённым
оператором» нельзя выполнить одновременно в положительном вещественном
представлении. Сохраняющий соотношение изоморфизм требует дополнительного
минуса и не является \(*\)-изоморфизмом принятой физической реализации.
\end{theorem}

Таким образом, удвоение колчана не устраняет выбор ориентации, а переносит
его в выбор симплектической поляризации между формальной обратной стрелкой
и эрмитовым сопряжением.

\section{Проверка вещественной калибровочной группы}

В действующей модели \(X,Y\) вещественны, поэтому
\begin{equation}
 \mu_G=XX^T-Y^TY
\end{equation}
симметрична. Для любой образующей \(\xi\in\mathfrak{so}(3)\), которая
кососимметрична,
\begin{equation}
 \Tr(\mu_G\xi)=0.
 \label{eq:v5-derived-so3-zero-pairing}
\end{equation}
Следовательно, \(\mu_G\) не является ненулевым строгим отображением момента
физической группы \(SO(3)\) при стандартном следовом отождествлении с
двойственной алгеброй Ли. Оно естественно для полного эрмитова или
матричного пространства, связанного с \(U(3)\) либо \(GL(3)\), но такое
замещение увеличивает размерность алгебры Ли с трёх до девяти и вводит шесть
новых образующих.

\begin{nogocriterion}[Скрытое расширение калибровочной группы]
Переход от \(SO(3)\) к \(U(3)\), \(GL(3)\) или их комплексной оболочке не
считается выводом из текущей модели, пока дополнительные образующие, поля и
их мера не получены из того же родительского объекта.
\end{nogocriterion}

\section{Что действительно даёт производный комплекс}

Минимальное разрешение Кошуля--Шафаревича добавляет образующие
\(\theta_i\) отрицательной степени и определяет
\begin{equation}
 \partial\theta_i=r_i,\qquad
 \partial x=\partial x^*=\partial y=\partial y^*=0.
 \label{eq:v5-derived-koszul-rule}
\end{equation}
Тогда \(\partial^2=0\) автоматически. Однако соотношения \(r_i\), включая
нужный знак в \(r_G\), уже стоят в правой части определения. Производное
разрешение правильно хранит их гомологию и факторизацию, но не выводит сами
соотношения из спектральных данных KO6.

Формализм Баталина--Вилковыского аналогично строит разрешение заданного
калибровочного действия и заданного главного функционала. Поэтому его
применение становится новым результатом только после независимого вывода
симплектической формы, отображения момента и группы, а не после записи
желаемого квадрата в главное уравнение.

\section{След и выбор средней вершины}

Идемпотент средней вершины \(e_G\) каноничен после задания размеченного
колчана. Он даёт точную нормировку
\begin{equation}
 \frac13\Tr\!\left(e_G[d,d^\dagger]^2\right)
 =\tau_3\!\left((XX^\dagger-Y^\dagger Y)^2\right).
 \label{eq:v5-derived-middle-trace}
\end{equation}
Но полный невзвешенный след содержит также квадраты \(X^\dagger X\) и
\(YY^\dagger\) на крайних вершинах. Чтобы оставить только
\eqref{eq:v5-derived-middle-trace}, необходимо заранее объявить среднюю
вершину калибруемой, а крайние --- обрамляющими. В проекте такая разметка
уже использовалась, но её родительское происхождение общей мерой не было
выведено.

\section{Граница с литературой}

Предпроективные алгебры и некоммутативная гамильтонова редукция дают
соотношение суммы коммутаторов и его независимость от первоначально
выбранной ориентации остова; см.
\path{arXiv:math/0502301} и \path{arXiv:2102.12336}.
Производные схемы представлений и производная пуассонова редукция строят
комплексы Кошуля и БРСТ для уже заданного отображения момента; см.
\path{arXiv:2006.09282} и \path{arXiv:2012.04451}.
Поэтому новый результат проекта состоит не в повторном получении известного
соотношения, а в выявлении несовместимости его формальной независимости от
ориентации с положительной \(*\)-реализацией и текущей вещественной группой.

\section{Вывод}

\begin{protocol}
\begin{enumerate}
 \item точное предпроективное соотношение в средней вершине:
 \textbf{получено};
 \item формальная независимость алгебры от направления стрелок:
 \textbf{выполнена};
 \item независимость при сохранении положительной инволюции:
 \textbf{не выполнена};
 \item ненулевое отображение момента физической группы \(SO(3)\):
 \textbf{не получено};
 \item нильпотентность производного комплекса:
 \textbf{выполнена условно после задания соотношения};
 \item вывод знака и группы из текущей геометрии KO6:
 \textbf{не выполнен};
 \item минимальная предпроективная ветвь как родительское объяснение:
 \textbf{закрыта};
 \item новая комплексная симплектическая теория:
 \textbf{остаётся возможной только как новая архитектура};
 \item физическое замыкание: \textbf{не достигнуто}.
\end{enumerate}
\end{protocol}

Следующая проверка ---
\path{version5_parent_architecture_status_freeze_gate}: необходимо свести
все ветви Тома V, отделить доказанные конструкции от условных и решить,
имеет ли смысл вводить новую комплексную симплектическую аксиоматику либо
следует заморозить программу на уровне строгих запретительных результатов.

Машинный сертификат сохранён в
\path{s2t_v5_derived_moment_map_minimal_data_gate_results.json}.